\section{Systemtests und Dauerbetrieb}
Nach erfolgreicher Integration aller Komponenten erfolgt die Prüfung des Gesamtsystems im Dauerbetrieb.  
Dabei werden Stabilität, Temperaturverhalten und Zuverlässigkeit untersucht.  

Das System wird zunächst über einen Zeitraum von mehreren Stunden betrieben.  
Währenddessen erfolgen wiederholte Kontrollen der WordClock-Anzeige.  
Die angezeigten Werte werden regelmäßig mit den entsprechenden Referenzquellen verglichen.  
Dabei kann keine Abweichung festgestellt werden.  
Zusätzlich wird überprüft, ob sich elektronische Komponenten oder Leitungen unzulässig erwärmen.  
Es treten keine kritischen Temperaturentwicklungen auf.  
Nachdem die WordClock mehrere Stunden in Betrieb war, wird diese noch mehrfach neu gestartet, um zu überprüfen, ob der Boot-Modus weiterhin zuverlässig funktioniert.  
Auch diese Überprüfung ist erfolgreich.  

Im nächsten Schritt erfolgt ein Dauerbetrieb über drei Tage.  
Dabei wird die Funktion der WordClock in regelmäßigen Abständen kontrolliert.  
Alle Anzeigen arbeiten während dieses Zeitraums stabil und fehlerfrei.  

Nach erfolgreichem Abschluss des Dauerbetriebs erfolgt der Aufbau der finalen Hardwareversion.  
Die zuvor verwendeten Breadboard- und Jumperleitungen werden durch fest verlötete Leitungen ersetzt.  
Zusätzlich werden die entwickelten und gedruckten Halterungen und Befestigungselemente montiert.  
Nachdem die Endmontage abgeschlossen ist, erfolgt eine erneute Überprüfung aller Funktionen und elektrischen Verbindungen.  
Dabei wird kontrolliert, ob alle Anzeigen korrekt arbeiten und sämtliche Leitungen ordnungsgemäß verbunden sind.  
Nach erfolgreichem Abschluss der Systemtests im Dauerbetrieb, wird die WordClock an ihrem vorgesehenen Einsatzort montiert und für die Abnahme in Betrieb genommen.